\section*{Chapitre \ref{ch:b2:muscle-movement} --- Contraction musculaire et motricité}

\begin{solution}{exo:b2:muscle-movement:1}
Les filaments fins glissent le long des filaments épais, tirés par les
têtes de myosine~; aucun des deux filaments ne raccourcit. La bande I et
la zone H se rétrécissent et les disques Z se rapprochent~; la bande A
(les filaments épais) garde sa longueur.
\end{solution}

\begin{solution}{exo:b2:muscle-movement:2}
(1) ATP fixé, tête détachée~; ATP hydrolysé, tête armée. (2) La tête se
fixe sur l'actine. (3) Libération du phosphate, coup de rame, libération
de l'ADP. (4) Un nouvel ATP se fixe et la tête se détache. L'ATP est
nécessaire pour \emph{libérer} la tête~: sans lui, toutes les têtes
restent fixées et le muscle se bloque --- c'est la rigidité.
\end{solution}

\begin{solution}{exo:b2:muscle-movement:3}
Influx nerveux~; libération d'acétylcholine et potentiel de plaque
motrice (\qty{0.6}{ms})~; \omterm{prop:b2:neurons-synapses:ap}{potentiel d'action} musculaire sur la fibre et
le long des \omterm{def:b2:cell-differentiation:fibre}{tubules T} (\qty{1}{ms})~; libération du calcium par le
\omterm{def:b2:cell-differentiation:fibre}{réticulum sarcoplasmique} (\qty{1}{ms})~; fixation du calcium sur la
troponine, déplacement de la tropomyosine, attachement des têtes
(quelques millisecondes)~; montée de la force en quelques dizaines de
millisecondes.
\end{solution}

\begin{solution}{exo:b2:muscle-movement:4}
\omterm{def:b2:muscle-movement:motorunit}{Unité motrice}~: un motoneurone et toutes les fibres qu'il innerve.
Secousse~: la contraction brève produite par un influx isolé.
\omterm{def:b2:muscle-movement:motorunit}{Tétanos}~: la contraction fusionnée et soutenue produite par des influx
qui arrivent plus vite que la secousse ne décroît. La force est graduée
par la fréquence de décharge de chaque unité et par le nombre d'unités
recrutées.
\end{solution}

\begin{solution}{exo:b2:muscle-movement:5}
Chaque \omterm{def:b2:cell-differentiation:fibre}{demi-sarcomère} glisse de \qty{0.2}{\micro m} en \qty{50}{ms}~:
\qty{4}{\micro m/s}~; cela représente 20 coups de rame de \qty{10}{nm} en
\qty{50}{ms}, soit 400 par seconde si une tête restait attachée tout du
long.
\end{solution}

\begin{solution}{exo:b2:muscle-movement:6}
\qty{2}{cm} en \qty{0.1}{s}~: \qty{0.2}{m/s}, 2 longueurs par seconde~;
chaque \omterm{def:b2:cell-differentiation:fibre}{sarcomère} raccourcit de $2/\num{40000}$ cm $= \qty{0.5}{\micro m}$ en
\qty{0.1}{s}, soit \qty{5}{\micro m/s}.
\end{solution}

\begin{solution}{exo:b2:muscle-movement:7}
$F/F_0 = 0.3125/(v/v_{\max} + 0.25) - 0.25$~: à 1, 3 et 6 longueurs par
seconde ($0.1$, $0.3$, $0.6$ de $v_{\max}$)~: $0.64$, $0.32$, $0.12$.
Puissance $Fv$ (en $F_0\times$ longueurs par seconde)~: $0.64$, $0.95$, $0.71$
--- maximale vers 3 longueurs par seconde, environ un tiers de $v_{\max}$.
\end{solution}

\begin{solution}{exo:b2:muscle-movement:8}
$80\times 30 = \qty{2400}{N}$~; au pied, $2400\times 5/45 =
\qty{270}{N}$.
\end{solution}

\begin{solution}{exo:b2:muscle-movement:9}
À \qty{10}{Hz}, les influx arrivent toutes les \qty{100}{ms}, une fois
disparu le calcium du précédent~: des secousses séparées, chacune partant
d'un état partiellement relâché, qui donnent une ondulation. À
\qty{50}{Hz}, ils arrivent toutes les \qty{20}{ms}, avant que le calcium
soit repompé~: le calcium reste élevé et la force fusionne. Le
\omterm{def:b2:muscle-movement:motorunit}{tétanos} dépasse la secousse parce qu'une impulsion calcique isolée est
trop brève pour que les ponts d'union tendent les éléments élastiques en
série et développent toute leur force~; un calcium soutenu le leur permet.
\end{solution}

\begin{solution}{exo:b2:muscle-movement:10}
Travail \qty{20}{kJ}, énergie de l'ATP \qty{80}{kJ}, soit $1.6$ mol
d'ATP. Pendant les \qty{5}{s} premières~: $0.8$ mol de phosphocréatine.
Les $0.8$ mol d'ATP restantes par glycolyse~: $0.4$ mol de glucose,
$0.8$ mol de lactate --- dans \qty{40}{L}, \qty{20}{mmol/L}.
\end{solution}

\begin{solution}{exo:b2:muscle-movement:11}
Canal de libération bloqué~: pas de calcium, pas de contraction ---
paralysie. Pompe bloquée~: le calcium reste élevé, le muscle ne peut pas
se relâcher --- une contracture. Troponine insensible~: la tropomyosine
ne se déplace jamais, pas de contraction. Pas de calcium dans le \omterm{def:b2:blood-circulation:blood}{sang}~: le
muscle squelettique se contracte normalement (son calcium est interne)~;
le \omterm{def:b2:heart:anatomy}{cœur}, qui a besoin du calcium extérieur pour déclencher sa libération,
s'arrête.
\end{solution}

\begin{solution}{exo:b2:muscle-movement:12}
La tête convertit en travail jusqu'à la moitié de l'énergie libre d'un
ATP~; le reste est de la chaleur, et les pertes des pompes, des ponts
d'union qui s'attachent sans tirer à grande vitesse et des éléments
élastiques ramènent l'ensemble à un quart --- le rendement d'un bon
moteur à essence. Deux manières de graduer la force permettent à la fois
un contrôle fin à faible effort (de petites unités ajoutées une à une,
chacune déchargeant plus ou moins vite) et une large gamme (un facteur
mille, d'une seule petite unité à la fréquence de secousse à toutes les
unités en \omterm{def:b2:muscle-movement:motorunit}{tétanos})~; un simple interrupteur ne donnerait qu'une force ou
aucune.
\end{solution}

\begin{solution}{pb:b2:muscle-movement:1}
\textbf{1.} $v = \sqrt{2\times 9.81\times 0.4} = \qty{2.8}{m/s}$~;
$\tfrac12\times 70\times 2.8^{2} = \qty{275}{J}$.
\textbf{2.} Travail $= 275 + 70\times 9.81\times 0.4 = \qty{550}{J}$ sur
\qty{0.4}{m}~: \qty{1370}{N}, deux fois le poids.
\textbf{3.} $F_0 = 500\times 30 = \qty{15000}{N}$~; au pied,
\qty{1500}{N}.
\textbf{4.} \qty{4}{cm} en \qty{0.3}{s}~: \qty{0.13}{m/s}, $1.1$ longueur
par seconde, $0.14\,v_{\max}$.
\textbf{5.} $F/F_0 = 0.3125/(0.14 + 0.25) - 0.25 = 0.55$~: \qty{8300}{N}
dans le muscle, \qty{830}{N} au pied --- bien en deçà des
\qty{1370}{N} nécessaires. Le raccourcissement musculaire seul ne peut
pas réaliser ce saut. Le contre-mouvement fournit le complément~: lorsque
le sauteur fléchit les jambes, les tendons étirés stockent de l'énergie
élastique, et le muscle, contracté de façon presque isométrique à force
élevée, les met en tension~; les tendons se détendent ensuite plus vite
qu'aucune fibre ne peut raccourcir.
\textbf{6.} $550/0.3 = \qty{1.8}{kW}$~; \qty{92}{W/kg} d'extenseur.
\textbf{7.} $4/\num{48000}$ cm $= \qty{0.83}{\micro m}$ par \omterm{def:b2:cell-differentiation:fibre}{sarcomère},
\qty{0.42}{\micro m} par \omterm{def:b2:cell-differentiation:fibre}{demi-sarcomère}~: 42 coups de rame de \qty{10}{nm}.
\textbf{8.} \qty{0.42}{\micro m} en \qty{0.3}{s}~: \qty{1.4}{\micro m/s}~;
140 coups de rame par seconde en cas d'attachement permanent.
\textbf{9.} $500\times 6\times 10^{10}\times \num{48000}\times 300 =
4.3\times 10^{20}$ têtes.
\textbf{10.} $0.55\times \num{15000}\times 0.04 = \qty{330}{J}$~; par
coup de rame, $0.4\times 8\times 10^{-20} = \qty{3.2e-20}{J}$~; $1.0\times
10^{22}$ coups de rame.
\textbf{11.} $1.0\times 10^{22}/4.3\times 10^{20} = 24$ coups de rame par
tête, sur les 42 disponibles~: environ \qty{60}{\%}. La réserve est
faible --- le muscle travaille près de sa limite, comme l'a montré la
question 5.
\textbf{12.} $1.0\times 10^{22}/6.0\times 10^{23} = 0.017$ mol,
\qty{8.7}{g}~; son énergie, $0.017\times 50 = \qty{860}{J}$, dont
\qty{330}{J} sont devenus du travail~: \qty{39}{\%}.
\textbf{13.} Plus d'un influx toutes les \qty{30}{ms}~: au-delà de
\qty{33}{Hz}~; environ 10 influx en \qty{0.3}{s}.
\textbf{14.} $9.9\times 10^{-6}\times 5\times 10^{-10}\times 6\times
10^{23} = 3\times 10^{9}$ ions libres~; $1.5\times 10^{9}$ ATP.
\textbf{15.} Fibres~: $500/(\pi\times(3\times 10^{-3})^{2}) = 1.8\times
10^{7}$~; ATP des ponts d'union par fibre et par influx~: $1.0\times
10^{22}/(1.8\times 10^{7}\times 10) = 6\times 10^{13}$, quarante mille
fois l'estimation du pompage. Le calcium libre n'est qu'un petit reliquat~:
le réticulum libère de l'ordre de \qty{0.1}{mmol/L}, dix fois
l'augmentation du calcium libre, presque entièrement fixé aussitôt par la
troponine et les tampons --- et il faut repomper la totalité.
\textbf{16.} Pour un saut maximal, presque toutes les unités en quelques
dizaines de millisecondes, les petites d'abord et les grandes en
dernier~; un pas tranquille en recrute peut-être un dixième à un cinquième.
\textbf{17.} Les unités sont recrutées par ordre de taille~: à faible
effort, chaque nouvelle unité ajoute un petit incrément~; à fort effort,
les unités restantes sont grandes et chacune ajoute un saut important.
\textbf{18.} Le potentiel de plaque motrice de chaque fibre est divisé
par deux~; les fibres dont le potentiel tombe sous le seuil ne déchargent
pas, si bien que la secousse et le \omterm{def:b2:muscle-movement:motorunit}{tétanos} sont plus faibles et que le
saut est moins haut.
\textbf{19.} $5\times 20 = \qty{0.1}{mol}$ d'ATP~: de quoi faire six sauts.
\textbf{20.} $20\times 20 = \qty{0.4}{mol}$ de phosphocréatine~:
environ 23 sauts.
\textbf{21.} Dix sauts utilisent $0.17$ mol, sans encore épuiser la
phosphocréatine~; après une vingtaine, la glycolyse prend le relais, et
le lactate et les protons s'accumulent.
\textbf{22.} \qty{1.8}{kW} de travail à \qty{40}{\%} représentent
\qty{4.6}{kW} d'ATP, $0.09$ mol par seconde, qui exigent $0.015$ mol
d'oxygène par seconde ---
\qty{0.34}{L/s}, \qty{21}{L/min}, dix fois ce que le \omterm{def:b2:blood-circulation:blood}{sang} peut apporter~:
le saut est anaérobie par nécessité, et l'oxygène vient après.
\textbf{23.} $0.17$ mol d'ATP exigent $0.029$ mol d'oxygène,
\qty{0.64}{L}~; remboursés en une minute, \qty{640}{mL/min}, deux fois et
demie la consommation de repos.
\textbf{24.} L'ATP est rephosphorylé à partir de la phosphocréatine en
quelques millisecondes, si bien que sa concentration baisse à peine~; la
rigidité exige un ATP presque nul, ce qu'un muscle vivant qui dispose de
la moindre réserve n'atteint jamais.
\textbf{25.} 42 coups de rame disponibles par \omterm{def:b2:cell-differentiation:fibre}{demi-sarcomère}, 24
nécessaires~; $F/F_0 = 0.55$ à la vitesse du saut~; $0.017$ mol,
\qty{8.7}{g}, d'ATP.
\end{solution}
